\documentclass[portrait, 11pt]{article}

\usepackage[portrait, margin=1in, top=0.75in, bottom=0.75in]{geometry}
\usepackage{xcolor, colortbl}
\usepackage{booktabs}
\usepackage{amssymb}
\usepackage{rotating}
\usepackage[most]{tcolorbox}
\usepackage{hyperref}
\usepackage{parskip}

\hypersetup{
    colorlinks=true,
    linkcolor=blue,
    filecolor=magenta,
    urlcolor=blue,
    citecolor=black,
    pdftitle={SNN Simulator Worksheet and Research Notes},
}

\urlstyle{same}

\title{SNN Simulator Worksheet and Research Notes}
\author{Lucas S Hindman \and Kurtis D Cantley}
\date{}

\begin{document}

\maketitle
\thispagestyle{empty}

\bigskip
\noindent
This document provides a capability summary worksheet for the open-source spiking neural
network (SNN) simulators surveyed in \emph{Finding Order in CHAOSS}, followed by
detailed research notes for each simulator. Each entry in the worksheet links to its
corresponding notes section via the \textbf{[notes]} reference. 

\clearpage

% -----------------------------------------------------------------------
% Worksheet (landscape fits naturally since the document is landscape)
% -----------------------------------------------------------------------

% Required Packages: xcolor, colortbl, boottabs, amssymb, rotating
\definecolor{HeaderGray}{gray}{0.70}
\definecolor{ColumnGray}{gray}{0.90}

\definecolor{SimGreen}{RGB}{208,230,202}
\definecolor{ModelYellow}{RGB}{254,239,193}
\definecolor{LearningOrange}{RGB}{251,223,194}
\definecolor{HardwareBlue}{RGB}{190,209,246}

\newcolumntype{a}{>{\columncolor{HeaderGray}}l}
\newcolumntype{e}{>{\columncolor{HeaderGray}}c}
\newcolumntype{d}{>{\columncolor{ColumnGray}}c}
\newcolumntype{f}{>{\columncolor{ModelYellow}}c}
\newcolumntype{g}{>{\columncolor{LearningOrange}}c}
\newcolumntype{h}{>{\columncolor{HardwareBlue}}c}
\newcolumntype{i}{>{\columncolor{SimGreen}}l}

\begin{table*}[htbp]
    \caption{Open-Source SNN Simulator Worksheet}
    \begin{center}
    \small
    \def\arraystretch{0.90}%  1 is the default, change whatever you need
    % Shades of Gray
    % \begin{tabular}{@{}|a||c|c|c|c|c||c|c|c||c|c|c|c||d|d|d||e|@{}}

    % Lots of Color
    \begin{tabular}{@{}|i||f|f|f|f|f||g|g|g||h|h|h|h||d|d|d||e|@{}}
        \toprule
        \rowcolor{HeaderGray}

        \multicolumn{1}{|c||}{\textbf{Simulator}} & \begin{turn}{90}\textbf{Bio-Inspired}\end{turn} & \begin{turn}{90}\textbf{Bio-Plausible}\end{turn} & \begin{turn}{90}\textbf{Integrate-and-Fire}\end{turn} & \begin{turn}{90}\textbf{Compartmental}\end{turn} & \begin{turn}{90}\textbf{Equation-Based}\end{turn} & \begin{turn}{90}\textbf{Supervised}\end{turn} & \begin{turn}{90}\textbf{Unsupervised}\end{turn} & \begin{turn}{90}\textbf{Static}\end{turn} & \begin{turn}{90}\textbf{CPU}\end{turn} & \begin{turn}{90}\textbf{GPU}\end{turn} & \begin{turn}{90}\textbf{FPGA}\end{turn} & \begin{turn}{90}\textbf{Neuromorphic}\end{turn} & \begin{turn}{90}\textbf{Freshness}\end{turn} & \begin{turn}{90}\textbf{Documentation}\end{turn} & \begin{turn}{90}\textbf{Community}\end{turn} & \begin{turn}{90}\textbf{Viability Score}\end{turn} \\
        \midrule
        \href{https://github.com/ANNarchy/ANNarchy}{ANNarchy}~\hyperref[ANNarchy]{[notes]} & \checkmark & \checkmark & \checkmark &   & \checkmark &   & \checkmark &   & \checkmark & \checkmark &   &   & 2 & 2 & 2 & 6 \\\hline
        \href{https://github.com/arbor-sim/arbor/}{Arbor}~\hyperref[Arbor]{[notes]} &   &   & \checkmark &   &   &   & \checkmark &   & \checkmark & \checkmark &   &   & 2 & 2 & 2 & 6 \\\hline
        \href{https://github.com/fzenke/auryn}{Auryn}~\hyperref[Auryn]{[notes]} &   &   & \checkmark &   &   &   & \checkmark &   & \checkmark &   &   &   & 2 & 2 & 1 & 5 \\\hline
        \href{https://github.com/BindsNET/bindsnet}{BindsNET}~\hyperref[BindsNET]{[notes]} & \checkmark &   & \checkmark &   &   &   & \checkmark & \checkmark & \checkmark & \checkmark &   &   & 2 & 2 & 2 & 6 \\\hline
        \href{https://github.com/BrainCog-X/Brain-Cog}{BrainCog}~\hyperref[BrainCog]{[notes]} & \checkmark & \checkmark & \checkmark & \checkmark &   & \checkmark & \checkmark &   & \checkmark & \checkmark & \checkmark &   & 2 & 2 & 2 & 6 \\\hline
        \href{https://github.com/brainpy/BrainPy}{BrainPy}~\hyperref[BrainPy]{[notes]} &   & \checkmark & \checkmark &   &   & \checkmark &   &   & \checkmark & \checkmark &   &   & 2 & 2 & 2 & 6 \\\hline
        \href{https://briansimulator.org/}{Brian2}~\hyperref[Brian2]{[notes]} &   &   &   & \checkmark & \checkmark &   & \checkmark &   & \checkmark & \checkmark &   &   & 2 & 2 & 2 & 6 \\\hline
        \href{https://github.com/sagacitysite/brian2_loihi}{Brian2Loihi}~\hyperref[Brian2Loihi]{[notes]} &   &   & \checkmark &   &   &   & \checkmark &   & \checkmark &   &   & \checkmark & 1 & 1 & 1 & 3 \\\hline
        \href{https://github.com/CRAFT-THU/BSim}{BSim}~\hyperref[BSim]{[notes]} &   &   & \checkmark &   &   &   &   & \checkmark & \checkmark & \checkmark &   &   & 0 & 0 & 1 & 1 \\\hline
        \href{https://github.com/UCI-CARL/CARLsim6}{CARLsim}~\hyperref[CARLsim]{[notes]} & \checkmark &   & \checkmark &   &   & \checkmark & \checkmark &   & \checkmark & \checkmark &   &   & 2 & 2 & 2 & 6 \\\hline
        \href{https://github.com/tudelft/cuSNN}{cuSNN}~\hyperref[cuSNN]{[notes]} &   &   & \checkmark &   &   &   & \checkmark & \checkmark &   & \checkmark &   &   & 0 & 1 & 1 & 2 \\\hline
        \href{https://github.com/helq/doryta}{Doryta}~\hyperref[Doryta]{[notes]} &   &   & \checkmark &   &   &   &   & \checkmark & \checkmark &   &   &   & 1 & 1 & 0 & 2 \\\hline
        \href{https://github.com/EduardoRosLab/edlut}{EDLUT}~\hyperref[EDLUT]{[notes]} & \checkmark & \checkmark & \checkmark &   &   &   & \checkmark & \checkmark & \checkmark & \checkmark &   &   & 2 & 0 & 1 & 3 \\\hline
        \href{https://github.com/emer/emergent}{Emergent}~\hyperref[Emergent]{[notes]} &   &   &   &   &   &   &   &   & \checkmark & \checkmark &   &   & 2 & 2 & 2 & 6 \\\hline
        \href{https://github.com/yhx/enlarge}{ENLARGE}~\hyperref[ENLARGE]{[notes]} &   &   & \checkmark &   &   &   &   & \checkmark &   & \checkmark &   &   & 1 & 1 & 1 & 3 \\\hline
        \href{https://github.com/sandialabs/Fugu}{Fugu}~\hyperref[Fugu]{[notes]} &   &   & \checkmark &   &   &   &   & \checkmark & \checkmark &   &   & \checkmark & 2 & 2 & 1 & 5 \\\hline
        \href{https://github.com/genesis-sim/genesis-2.4}{Genesis}~\hyperref[Genesis]{[notes]} &   & \checkmark &   & \checkmark &   &   & \checkmark &   & \checkmark &   &   &   & 0 & 2 & 1 & 3 \\\hline
        \href{https://github.com/genn-team/genn}{GeNN}~\hyperref[GeNN]{[notes]} & \checkmark & \checkmark &   &   & \checkmark &   & \checkmark &   & \checkmark & \checkmark &   &   & 2 & 2 & 2 & 6 \\\hline
        \href{https://github.com/electronicvisions/jaxsnn}{jaxSNN}~\hyperref[jaxSNN]{[notes]} &   &   & \checkmark &   &   &   &   & \checkmark & \checkmark &   &   & \checkmark & 2 & 2 & 1 & 5 \\\hline
        \href{https://github.com/lava-nc/lava}{Lava}~\hyperref[Lava]{[notes]} &   &   & \checkmark &   &   &   & \checkmark &   & \checkmark &   &   & \checkmark & 2 & 2 & 2 & 6 \\\hline
        \href{https://github.com/YLab-UChicago/Memristor-Spikelearn}{Memristor-Spikelearn}~\hyperref[Memristor-Spikelearn]{[notes]} &   &   & \checkmark &   &   &   & \checkmark &   & \checkmark &   &   &   & 2 & 1 & 0 & 3 \\\hline
        \href{https://git.renater.fr/anonscm/git/n2s3/n2s3.git/}{N2S3}~\hyperref[N2S3]{[notes]} &   &   & \checkmark &   &   &   & \checkmark &   & \checkmark &   &   & \checkmark & 0 & 2 & 2 & 4 \\\hline
        \href{https://www.nengo.ai/}{Nengo}~\hyperref[Nengo]{[notes]} & \checkmark &   & \checkmark &   &   &   & \checkmark &   & \checkmark & \checkmark & \checkmark & \checkmark & 2 & 2 & 2 & 6 \\\hline
        \href{https://github.com/nengo/nengo-dl}{NengoDL}~\hyperref[NengoDL]{[notes]} & \checkmark &   & \checkmark &   &   & \checkmark &   & \checkmark & \checkmark & \checkmark &   &   & 2 & 2 & 2 & 6 \\\hline
        \href{https://github.com/nest/nest-simulator}{NEST}~\hyperref[NEST]{[notes]} & \checkmark & \checkmark & \checkmark & \checkmark & \checkmark &   & \checkmark & \checkmark & \checkmark &   &   &   & 2 & 2 & 2 & 6 \\\hline
        \href{https://github.com/nest/nest-gpu}{NEST GPU}~\hyperref[NEST GPU]{[notes]} &   & \checkmark & \checkmark &   & \checkmark &   & \checkmark &   &   & \checkmark &   &   & 2 & 2 & 1 & 5 \\\hline
        \href{https://github.com/neuronsimulator/nrn}{Neuron}~\hyperref[Neuron]{[notes]} &   & \checkmark & \checkmark & \checkmark &   &   & \checkmark &   & \checkmark &   &   &   & 2 & 2 & 2 & 6 \\\hline
        \href{https://github.com/hjq310/NeuroPack}{NeuroPack}~\hyperref[NeuroPack]{[notes]} & \checkmark &   & \checkmark &   &   & \checkmark & \checkmark &   & \checkmark &   &   &   & 1 & 1 & 0 & 2 \\\hline
        \href{https://github.com/hoangt/neutrams}{NEUTRAMS}~\hyperref[NEUTRAMS]{[notes]} &   &   & \checkmark &   &   &   &   & \checkmark & \checkmark &   &   & \checkmark & 0 & 0 & 0 & 0 \\\hline
        \href{https://github.com/norse/norse}{Norse}~\hyperref[Norse]{[notes]} & \checkmark &   & \checkmark &   &   &   & \checkmark &   & \checkmark & \checkmark &   & \checkmark & 2 & 2 & 2 & 6 \\\hline
        \href{https://github.com/trieschlab/PymoNNto}{PymoNNto}~\hyperref[PymoNNto]{[notes]} & \checkmark & \checkmark & \checkmark &   &   & \checkmark & \checkmark &   & \checkmark & \checkmark &   &   & 2 & 2 & 1 & 5 \\\hline
        \href{https://github.com/cnrl/PymoNNtorch}{PymoNNtorch}~\hyperref[PymoNNtorch]{[notes]} & \checkmark & \checkmark & \checkmark &   &   & \checkmark & \checkmark &   & \checkmark & \checkmark &   &   & 2 & 2 & 1 & 5 \\\hline
        \href{https://github.com/BasBuller/PySNN}{PySNN}~\hyperref[PySNN]{[notes]} &   &   & \checkmark &   &   &   & \checkmark &   & \checkmark & \checkmark &   &   & 0 & 2 & 1 & 3 \\\hline
        \href{https://github.com/UA-RCL/RANC}{RANC}~\hyperref[RANC]{[notes]} &   &   & \checkmark &   &   & \checkmark &   &   & \checkmark &   & \checkmark &   & 1 & 0 & 1 & 2 \\\hline
        \href{https://github.com/Rao-Sanaullah/RAVSim}{RAVSim}~\hyperref[RAVSim]{[notes]} &   &   & \checkmark &   &   &   &   &   & \checkmark &   &   &   & 2 & 2 & 0 & 4 \\\hline
        \href{https://github.com/WISION-Lab/sarnn}{SaRNN}~\hyperref[SaRNN]{[notes]} &   &   & \checkmark &   &   &   &   & \checkmark & \checkmark &   &   & \checkmark & 2 & 0 & 0 & 2 \\\hline
        \href{https://github.com/synsense/sinabs}{SINABS}~\hyperref[SINABS]{[notes]} &   &   & \checkmark &   &   & \checkmark &   & \checkmark & \checkmark &   &   & \checkmark & 2 & 2 & 2 & 6 \\\hline
        \href{https://github.com/NeuromorphicProcessorProject/snn_toolbox}{SNNToolbox}~\hyperref[SNNToolbox]{[notes]} &   &   & \checkmark &   &   &   &   & \checkmark & \checkmark &   &   &   & 1 & 2 & 2 & 5 \\\hline
        \href{https://github.com/jeshraghian/snntorch}{snnTorch}~\hyperref[snnTorch]{[notes]} &   &   & \checkmark &   &   & \checkmark &   &   & \checkmark & \checkmark &   & \checkmark & 2 & 2 & 2 & 6 \\\hline
        \href{https://github.com/zju-bmi-lab/SPAIC}{SPAIC}~\hyperref[SPAIC]{[notes]} & \checkmark & \checkmark & \checkmark &   &   & \checkmark & \checkmark &   & \checkmark & \checkmark &   &   & 2 & 2 & 1 & 5 \\\hline
        \href{https://github.com/OFTNAI/Spike}{SPIKE}~\hyperref[SPIKE]{[notes]} &   &   & \checkmark &   &   & \checkmark &   & \checkmark &   & \checkmark &   &   & 0 & 2 & 1 & 3 \\\hline
        \href{https://github.com/fangwei123456/spikingjelly}{SpikingJelly}~\hyperref[SpikingJelly]{[notes]} &   &   & \checkmark &   &   & \checkmark & \checkmark & \checkmark & \checkmark & \checkmark &   & \checkmark & 2 & 2 & 2 & 6 \\\hline
        \href{https://github.com/miladmozafari/SpykeTorch}{SpykeTorch}~\hyperref[SpykeTorch]{[notes]} &   &   & \checkmark &   &   &   & \checkmark &   & \checkmark & \checkmark &   &   & 1 & 2 & 0 & 3 \\\hline
        \href{https://github.com/kmheckel/spyx}{Spyx}~\hyperref[Spyx]{[notes]} &   &   & \checkmark &   &   & \checkmark &   &   &   & \checkmark &   & \checkmark & 2 & 2 & 1 & 5 \\\hline
        \href{https://github.com/ORNL/superneuro}{SuperNeuro}~\hyperref[SuperNeuro]{[notes]} &   &   & \checkmark &   &   &   & \checkmark &   & \checkmark & \checkmark &   &   & 2 & 1 & 1 & 4 \\\hline

        \bottomrule
        \rowcolor{HeaderGray}
        &  \multicolumn{5}{c}{\textbf{Neuron Model}} & \multicolumn{3}{c}{\textbf{Learning}} & \multicolumn{4}{c}{\textbf{Hardware}} & \multicolumn{4}{c}{\textbf{Viability}} \\
        \rowcolor{HeaderGray}
        &  \multicolumn{5}{c}{\textbf{Family}} & \multicolumn{3}{c}{\textbf{Method}} & \multicolumn{4}{c}{\textbf{Support}} & \multicolumn{4}{c}{\textbf{Scorecard}} \\
    \end{tabular}
    \label{simWorksheet}
    \end{center}
\end{table*}


\clearpage

% -----------------------------------------------------------------------
% Per-simulator research notes
% -----------------------------------------------------------------------
\section*{Simulator Research Notes}

\input{simulator-notes}

\end{document}
